\chapter{Childhood obesity}
\index{Childhood obesity}

\medskip
\noindent\textit{Educational information; seek professional advice for individual care.}

\section*{What it means}
Childhood obesity means excess body fat that may affect health. It is influenced by genetics, appetite, sleep, medicines, food access, activity, stress, and the wider environment; it is not simply a matter of willpower.

\section*{How it is assessed}
Health professionals track growth over time using age- and sex-specific growth charts. For children under 5, weight-for-height is commonly used; for ages 5--19, body-mass-index-for-age may be used. These measures are screening tools, not a complete judgement of health. Assessment may also include blood pressure, sleep, emotional wellbeing, medicines, eating patterns, activity, and signs of related conditions.

\section*{What can help}
The aim is healthy growth and wellbeing, not rapid weight loss. Families can work with a clinician or dietitian on regular meals, water as the usual drink, varied foods, enjoyable activity, adequate sleep, and less sedentary screen time. Changes should be realistic, family-wide, and free from shame. Do not put a child on a restrictive diet or use weight-loss medicines or supplements without specialist advice.

\section*{When to seek medical advice}
Arrange a review if weight is changing rapidly, growth is unusual, the child snores or pauses breathing during sleep, has exercise intolerance, joint pain, excessive thirst or urination, low mood, eating concerns, or bullying. Seek urgent care for severe breathing difficulty, collapse, confusion, or other rapidly worsening illness.

\section*{Follow-up}
Follow-up depends on age, growth pattern, health risks, and family goals. Treat related conditions and review progress without relying on a single measurement. Ask for support if food insecurity, disability, medicines, or mental-health concerns make changes difficult.

\section*{Sources}
\begin{itemize}
\item World Health Organization: \url{https://www.who.int/news-room/questions-and-answers/item/noncommunicable-diseases-childhood-overweight-and-obesity}
\item World Health Organization child growth standards: \url{https://www.who.int/tools/child-growth-standards}
\end{itemize}
